% Drawing library for the craft sheets.
%
% The craft sheets teach transferable shop skill, not launcher assembly, so the
% drawings are of tools and joints rather than of this project's parts. A boy
% who learns to cut square and solvent-weld properly here should recognise the
% same drawings when he builds something else entirely.
%
% Everything is drawn in section where a section shows the thing that matters.
% A solvent weld is invisible from outside, and that is exactly why it gets
% built wrong.

\tikzset{
  craftline/.style={draw=telosink,line width=0.9pt},
  craftheavy/.style={draw=telosink,line width=1.3pt},
  craftdetail/.style={draw=telosink,line width=0.35pt},
  crafthint/.style={draw=telosmute,line width=0.4pt,dash pattern=on 2pt off 1.6pt},
  craftfill/.style={fill=telosfill},
  craftweld/.style={pattern=north east lines,pattern color=telosdeep},
  craftbad/.style={draw=telosink,line width=0.9pt},
  craftpencil/.style={draw=telosink,line width=0.28pt,opacity=0.72},
  craftgrain/.style={draw=telosmute,line width=0.22pt,opacity=0.62},
  cflabel/.style={font=\fontsize{6.4}{7.4}\selectfont},
  cfsmall/.style={font=\fontsize{5.8}{6.8}\selectfont},
  cftag/.style={font=\fontsize{6.4}{7.4}\selectfont\itshape},
}

% ---------------------------------------------------------------------------
% Pipe and fittings, in longitudinal section
% ---------------------------------------------------------------------------

% A length of pipe in section. \pipesection{x}{y}{length}{outer radius}{wall}
\newcommand{\pipesection}[5]{%
  \begin{scope}[shift={(#1,#2)}]
    \draw[craftline,craftfill] (0,#4-#5) rectangle (#3,#4);
    \draw[craftline,craftfill] (0,-#4) rectangle (#3,-#4+#5);
    \draw[crafthint] (0,0) -- (#3,0);
  \end{scope}}

% A socket fitting in section, with the pipe inserted to the shoulder and the
% interference zone marked. This is the drawing that explains solvent welding.
%   \socketjoint{x}{y}{scale}{insert depth fraction 0..1}
\newcommand{\socketjoint}[4]{%
  \begin{scope}[shift={(#1,#2)},scale=#3]
    % fitting body
    \draw[craftline,craftfill] (1.4,0.30) rectangle (3.4,0.62);
    \draw[craftline,craftfill] (1.4,-0.30) rectangle (3.4,-0.62);
    \draw[craftline] (3.4,0.62) -- (3.4,-0.62);
    % socket shoulder
    \draw[craftline] (2.9,0.30) -- (2.9,-0.30);
    % pipe, inserted #4 of the way to the shoulder
    \pgfmathsetmacro{\tip}{1.4 + (2.9-1.4)*#4}
    \draw[craftline,craftfill] (-0.1,0.30) rectangle (\tip,0.44);
    \draw[craftline,craftfill] (-0.1,-0.30) rectangle (\tip,-0.44);
    % the weld zone: where softened surfaces fuse
    \fill[craftweld] (1.4,0.30) rectangle (\tip,0.44);
    \fill[craftweld] (1.4,-0.30) rectangle (\tip,-0.44);
    \draw[crafthint] (-0.1,0) -- (3.4,0);
  \end{scope}}

% A mitre box holding pipe, seen from above.
\newcommand{\mitrebox}[3]{%
  \begin{scope}[shift={(#1,#2)},scale=#3]
    % Bench-top perspective, with enough construction detail to read as a tool.
    \fill[craftfill] (0,0) -- (3.2,0) -- (3.48,0.20) -- (0.28,0.20) -- cycle;
    \draw[craftheavy] (0,0) -- (3.2,0) -- (3.48,0.20) -- (3.48,1.42)
      -- (3.2,1.50) -- (0,1.50) -- (0.28,1.42) -- (0.28,0.20) -- cycle;
    \draw[craftline] (0,0.28) -- (3.2,0.28);
    \draw[craftline] (0,1.22) -- (3.2,1.22);
    \draw[craftdetail] (3.2,0) -- (3.2,1.50);
    \foreach \y in {0.38,0.52,...,1.12} {
      \draw[craftgrain] (0.08,\y) .. controls (0.8,\y+0.04) and
        (2.35,\y-0.05) .. (3.12,\y+0.01);}
    % the 90-degree slot
    \draw[craftheavy] (1.5,-0.12) -- (1.5,1.62);
    \draw[craftheavy] (1.62,-0.12) -- (1.62,1.62);
    % pipe lying in the channel
    \draw[craftline,fill=telospaper] (0.2,0.42) rectangle (2.9,1.08);
    \draw[craftdetail] (0.2,0.42) -- (0.2,1.08);
    \draw[craftpencil] (0.2,0.50) -- (2.9,0.50);
    \draw[craftpencil] (0.2,1.00) -- (2.9,1.00);
    \draw[craftdetail] (0.2,0.42) arc[start angle=-90,end angle=90,
      x radius=0.13,y radius=0.33];
    \draw[craftpencil] (1.44,0.41) -- (1.44,1.09);
    \draw[craftpencil] (1.68,0.41) -- (1.68,1.09);
  \end{scope}}

% A try square checking a cut end.
\newcommand{\squarecheck}[4]{%
  \begin{scope}[shift={(#1,#2)},scale=#3]
    % pipe end, cut at angle #4 degrees off square
    \draw[craftline,craftfill] (0,0.28) rectangle (1.6,0.46);
    \draw[craftline,craftfill] (0,-0.28) rectangle (1.6,-0.46);
    \pgfmathsetmacro{\skewa}{1.6 + 0.46*tan(#4)}
    \pgfmathsetmacro{\skewb}{1.6 - 0.46*tan(#4)}
    \draw[craftline] (\skewa,0.46) -- (\skewb,-0.46);
    % square: stock against the pipe side, blade across the end
    \draw[craftheavy] (1.1,0.46) -- (1.1,0.86) -- (2.5,0.86);
    \draw[craftheavy] (1.1,0.46) -- (2.35,0.46);
    \draw[craftdetail] (1.24,0.46) -- (1.24,0.72);
    \foreach \x in {1.35,1.55,...,2.35} {
      \draw[craftpencil] (\x,0.86) -- (\x,0.78);}
    \fill[telospaper] (1.02,0.53) rectangle (1.30,0.76);
    \draw[craftline] (1.02,0.53) rectangle (1.30,0.76);
    \draw[craftdetail] (1.08,0.59) -- (1.24,0.59) -- (1.24,0.70);
  \end{scope}}

% A chamfered pipe end in section, with the angle called out.
\newcommand{\chamferend}[3]{%
  \begin{scope}[shift={(#1,#2)},scale=#3]
    \draw[craftline,craftfill] (0,0.30) -- (1.5,0.30) -- (1.5,0.44)
      -- (0.22,0.44) -- (0,0.30) -- cycle;
    \draw[craftline,craftfill] (0,-0.30) -- (1.5,-0.30) -- (1.5,-0.44)
      -- (0.22,-0.44) -- (0,-0.30) -- cycle;
    \draw[crafthint] (0,0) -- (1.5,0);
    \draw[crafthint] (0.22,0.44) -- (0.42,0.44);
    \draw[crafthint] (0,0.30) -- (0.42,0.30);
    \foreach \x in {0.25,0.38,...,1.42} {
      \draw[craftgrain] (\x,0.32) -- (\x+0.10,0.42);
      \draw[craftgrain] (\x,-0.32) -- (\x+0.10,-0.42);}
  \end{scope}}

% A burr on an inside edge -- the defect the chamfer sheet is about.
\newcommand{\burrend}[3]{%
  \begin{scope}[shift={(#1,#2)},scale=#3]
    \draw[craftline,craftfill] (0,0.30) rectangle (1.5,0.44);
    \draw[craftline,craftfill] (0,-0.30) rectangle (1.5,-0.44);
    \draw[craftline] (0,0.30) .. controls (0.06,0.20) and (0.14,0.18) .. (0.10,0.30);
    \draw[craftline] (0,-0.30) .. controls (0.06,-0.20) and (0.14,-0.18) .. (0.10,-0.30);
    \draw[crafthint] (0,0) -- (1.5,0);
    \foreach \x in {0.08,0.22,...,1.42} {
      \draw[craftgrain] (\x,0.32) -- (\x+0.10,0.42);
      \draw[craftgrain] (\x,-0.32) -- (\x+0.10,-0.42);}
    \draw[craftpencil] (0.02,0.29) .. controls (0.12,0.17) and
      (0.20,0.22) .. (0.15,0.31);
  \end{scope}}

% A tick-and-cross verdict marker, for right-way/wrong-way pairs.
\newcommand{\craftyes}[2]{%
  \draw[craftheavy] (#1-0.14,#2) -- (#1-0.04,#2-0.12) -- (#1+0.16,#2+0.14);}
\newcommand{\craftno}[2]{%
  \draw[craftheavy] (#1-0.13,#2-0.13) -- (#1+0.13,#2+0.13);
  \draw[craftheavy] (#1-0.13,#2+0.13) -- (#1+0.13,#2-0.13);}

% Dimension line with a label.
\newcommand{\craftdim}[5]{%
  \draw[|<->|,craftdetail] (#1,#2) -- (#3,#4);
  \node[cftag,fill=telospaper,inner sep=1.6pt] at ({(#1+#3)/2},{(#2+#4)/2}) {#5};}

% Callout with a hairline leader.
\newcommand{\craftcall}[6]{%
  \draw[craftdetail,telosmute] (#1,#2) -- (#3,#4);
  \fill[telosmute] (#1,#2) circle (0.035);
  \node[cflabel,anchor=#5,align=left,inner sep=1.6pt] at (#3,#4) {#6};}

% ---------------------------------------------------------------------------
% Craft-sheet furniture
% ---------------------------------------------------------------------------

% The standard every craft sheet opens with: what "good" looks like, stated as
% something checkable rather than as an aspiration.
\newtcolorbox{standard}{
  enhanced, breakable,
  colback=telospaper, colframe=telosink,
  boxrule=1.0pt, sharp corners,
  left=6pt, right=6pt, top=4pt, bottom=4pt,
  toptitle=3pt, bottomtitle=3pt, titlerule=0.4pt,
  before skip=7pt, after skip=7pt,
  colbacktitle=telospaper, coltitle=telosink,
  title={THE STANDARD}, fonttitle=\small\bfseries
}

% What you should observe if the step worked.
\newcommand{\expectcraft}[1]{%
  \par\vspace{0.15em}
  \noindent\begin{minipage}{\linewidth}
  \hspace*{0.6em}\begin{minipage}{\dimexpr\linewidth-0.6em\relax}
  {\footnotesize\textbf{You should see:} #1\par}
  \end{minipage}\end{minipage}\par\vspace{0.3em}}

% How to check your own work. Every craft sheet has one, because a skill you
% cannot verify is a skill you cannot trust.
\newcommand{\checkyourself}[1]{%
  \par\vspace{0.3em}
  \noindent
  \begin{minipage}{\linewidth}
    \rule{\linewidth}{0.9pt}\par
    \vspace{0.3em}
    {\footnotesize\textbf{Check it yourself.} #1\par}
    \vspace{0.25em}
    \rule{\linewidth}{0.3pt}
  \end{minipage}\par}
